\documentclass[12pt,a4paper]{article}
\usepackage[utf8]{inputenc}
\usepackage{geometry}
\geometry{margin=1in}
\usepackage{graphicx}
\usepackage{titlesec}
\usepackage{tocloft}
\usepackage{hyperref}
\usepackage{fancyhdr}  % For headers and footers
\hypersetup{
    colorlinks=true,
    linkcolor=blue,
    urlcolor=blue,
}

% Customize section formatting
\titleformat{\section}{\large\bfseries}{\thesection}{1em}{}
\titleformat{\subsection}{\normalsize\bfseries}{\thesubsection}{1em}{}

% Header and footer setup
\pagestyle{fancy}
\fancyhf{}  % Clear all headers and footers

% Left header
\fancyhead[L]{\textbf{Scope Review}}

% Right header
\fancyhead[R]{\textbf{LV Isolated DC/DC Converter}}

% Page number in footer
\fancyfoot[C]{\thepage}

% Ensure the header appears on all pages, including the first
\thispagestyle{fancy}

\begin{document}

% Title Page
\begin{center}
    {\huge \textbf{LV Isolated DC/DC Converter}} \\[0.3cm]
    {\large AUTO1931 --- Scope Report} \\[0.2cm]
    {\today}\\[0.5cm]

    {\large William Bowley}
    \vfill
\end{center}

% Foreword
\section*{Foreword}
\addcontentsline{toc}{section}{Foreword}
This scope report establishes the architectural design of the Isolated DC/DC converter and defines the implementation boundaries of the system. It is presented as a series of 
conceptual steps: from system decomposition and requirement discovery, through the understanding and modeling of a selected topology, to the conceptual design that emerged from
that process. The document concludes with project management considerations and general conclusions prior to implementation.

% Table of Contents
\tableofcontents
\newpage

% ========== BACKGROUND ==========
\section{Background}\label{sec:background}
An isolated DC/DC converter is a component within the tractive and low-voltage grounded (LVG) system. Its purpose is to convert the tractive battery voltage (approximately 500–600 V) 
to a regulated 12 V with minimal ripple, while maintaining galvanic isolation. This component also charges the LVG battery, which serves as an energy buffer and provides standby power when
 the tractive battery is removed from the vehicle.

\vspace{0.5cm}

The DC/DC converter removes the constraint of LVG battery capacity, improving endurance while potentially reducing packaging. Furthermore, according to Archer et al. (2024b), 
it eliminates challenges associated with separate charging arrangements. This system also potentially enables future expansion of the LVG system to support actuators for driverless use cases.

% ========== SYSTEM SPECIFICATION ==========
\section{System Specification}\label{sec:specification}




Here, you are to specify what your system is, the targets of this system and identify the major system design constraints. 
This section is to serve as the effective Problem Statement for your project. As such, it must be clear and obvious what the objectives of your system are.

% ========== LITERATURE AND DESIGN REVIEW ==========
\section{Literature and Design Review}\label{sec:litreview}
With your problem developed, you may now present, analyse and interpret external literature related to your project. 
It is anticipated that all presented literature will assist in informing the overall and specific details of your design work --- if the literature does not inform your work, do not include that literature.

\subsection{Example of Literature}\label{subsec:example}
Your literature and design review may include:
\begin{itemize}
    \item Design reports from former RMIT FSAE teams and other Formula SAE teams
    \item Technical reports
    \item Key information from relevant textbooks
    \item High quality journal or technical papers from reputable sources
    \begin{itemize}
        \item Do not review reports from low quality sources (e.g.\@ pay-to-publish academic journals).
    \end{itemize}
    \item Reviews of existing eminent --- or otherwise --- designs relevant to your project
    \begin{itemize}
        \item This can include social media posts and team or personal photography
    \end{itemize}
    \item Review of relevant production techniques, including materials, manufacturing process and off-the-shelf components.
\end{itemize}

\subsection{Interpretation}\label{subsec:interpretation}
Do not simply present the literature. You are to rigorously review the literature in the context of your problem. 
It is incumbent on you --- as the design engineer --- to interpret the literature and communicate the influence of that literature on your project. 
This may include, and is definitively not limited to, the critical analysis of designs from other Formula vehicles, evaluation of the Overall Vehicle 
Specification on your system or concise summaries of engineering science and theory pertaining to your system.

% ========== CONCEPTUAL DESIGN ==========
\section{Conceptual Design}\label{sec:conceptual}
At this stage of the report, it should be clear what the requirements of your system are 
(from System Specification section) and what the engineering context of your design problem is 
(from the Literature and Design Review).

Now you are to present and discuss the overall concept of your system. You must indicate how 
the features of the concept intend to meet the targets of the system, in a general manner. 
You may make use of appendices to communicate information or present figures which do not fit into the page limitations of the report.

% ========== PROJECT MANAGEMENT ==========
\section{Project Management}\label{sec:project}

Your Project Management plan should consist of a definition of project deliverables, an analysis of required resources (which may include materials,
knowledge and data from RMIT's FSAE teams) and a timeline (preferably Gantt Chart) for the project. Please ensure that the timeline is effectively 
in visually communicating when you intend for events to occur.

The focus of this section should be regarding the operation and completion of the design cycle.

% ========== CONCLUSION ==========
\section{Conclusion}\label{sec:conclusion}
Summarise the outcomes of the report. Nominate any areas where particular progress or novel outcomes have been achieved.

\end{document}